\ifx\allfiles\undefined
\documentclass[12pt, a4paper, oneside, UTF8]{ctexbook}
\def\path{../config}
\input{\path/package.tex}

% 定理环境设置
\input{\path/theorem1_zh.tex}

\input{\path/custom.tex}

% 修改参数改变封面样式，0 默认原始封面、内置其他1、2、3种封面样式
\def\myIndex{3}


\ifnum\myIndex>0
    \input{\path/cover_package_\myIndex} 
\fi

\def\myTitle{复变函数与积分变换笔记}
\def\myAuthor{M.K.}
\def\myDateCover{\today}
\def\myDateForeword{\today}
\def\myForeword{前言}
\def\myForewordText{
    
    这是一个基于\LaTeX{}模板的复变函数与积分变换笔记．
    内容参考了包括但不限于以下资料：
    \begin{itemize}
        \item 《复变函数与积分变换》（科学出版社）第三版
        \item 《复变函数习题精选精讲》（山东科学技术出版社）第一版
    \end{itemize}

    封面来源于\href{https://www.pixiv.net/users/2642047}{アシマ / Ashima}的作品\href{https://www.pixiv.net/artworks/147092236}{Assembly}
    如有侵权请联系我删除．
}
\def\mySubheading{副标题}


\begin{document}
%\input{../config/cover}
\else
\fi
\chapter{复积分}

\section{复积分}

\begin{definition}[逐段光滑曲线]
设 $\gamma\colon [a,b]\to\C$，$\gamma(t)=x(t)+iy(t)$。若 $x(t),y(t)$ 在 $[a,b]$ 上连续且除有限个点外在 $[a,b]$ 上连续可微，则称 $\gamma$ 为\textbf{逐段光滑曲线}。
\end{definition}

\begin{definition}[复积分]
设 $\gamma\colon [a,b]\to\C$ 为逐段光滑曲线，$f$ 在 $\gamma$ 上有界。若极限
\[
\int_\gamma f(z)\,\d z \defeq \lim_{n\to\infty}\sum_{k=1}^{n}f(\zeta_k)(z_k-z_{k-1})
\]
存在（与分割 $\{z_k\}$ 及介点 $\{\zeta_k\}$ 的选取无关），则称 $f$ 沿 $\gamma$\textbf{可积}，称该极限为 $f$ 沿 $\gamma$ 的\textbf{复积分}（或\textbf{第一型复积分}）。
\end{definition}

\begin{theorem}[复积分的等价形式]
设 $f(z)=u(x,y)+iv(x,y)$，$\gamma$ 逐段光滑。则 $f$ 沿 $\gamma$ 可积当且仅当 $u,v$ 沿 $\gamma$ 可积，且
\[
\int_\gamma f(z)\,\d z = \int_\gamma u\,\d x - v\,\d y + i\int_\gamma v\,\d x + u\,\d y.
\]
\end{theorem}

\begin{theorem}[积分基本性质]
设 $f,g$ 沿逐段光滑曲线 $\gamma$ 可积，$\alpha,\beta\in\C$，则
\begin{enumerate}
    \item 线性性：$\displaystyle\int_\gamma (\alpha f + \beta g)\,\d z = \alpha\int_\gamma f\,\d z + \beta\int_\gamma g\,\d z$；
    \item 可加性：若 $\gamma=\gamma_1+\gamma_2$，则 $\displaystyle\int_\gamma f\,\d z = \int_{\gamma_1}f\,\d z + \int_{\gamma_2}f\,\d z$；
    \item 反向性：$\displaystyle\int_{\gamma^-} f\,\d z = -\int_\gamma f\,\d z$，其中 $\gamma^-$ 表示 $\gamma$ 的反向曲线；
    \item 绝对值估计：$\displaystyle\abs{\int_\gamma f\,\d z} \le \max_{z\in\gamma}\abs{f(z)}\cdot \int_{\gamma} \abs{\gamma'(t)}\,\d t$；
    \item 共轭：$\displaystyle\overline{\int_\gamma f\,\d z} = \int_\gamma \overline{f(z)}\,\d\bar z$。
\end{enumerate}
\end{theorem}

\begin{example}
    求积分 $\displaystyle\int_C\frac{\d z}{(z-z_0)^n},\;n\in \N$，其中 $C$ 为以 $z_0$ 为圆心、半径为 $R>0$ 的圆，取逆时针方向。
\end{example}
\begin{solution}

    设 $z = z_0 + Re^{i\theta}$，则 $\d z = iRe^{i\theta}\d\theta$，$\theta\in[0,2\pi]$。于是
    \[
    \int_C\frac{\d z}{(z-z_0)^n} = \int_0^{2\pi}\frac{iRe^{i\theta}\d\theta}{(Re^{i\theta})^n} = iR^{1-n}\int_0^{2\pi}e^{i(1-n)\theta}\d\theta.
    \]
    \begin{enumerate}
        \item 当 $n=1$ 时，$\displaystyle\int_0^{2\pi}e^{i(1-n)\theta}\d\theta = \int_0^{2\pi}\d\theta = 2\pi$
        \item 当 $n>1$ 时，$\displaystyle\int_0^{2\pi}e^{i(1-n)\theta}\d\theta = \left[\frac{e^{i(1-n)\theta}}{i(1-n)}\right]_0^{2\pi} = 0$
    \end{enumerate}

    因此
    \[
    \int_C\frac{\d z}{(z-z_0)^n} =
    \begin{cases}
        2\pi i, & n=1\\
        0, & n>1.
    \end{cases}
    \]
\end{solution}
\section{柯西积分定理}

\begin{theorem}[Cauchy--Goursat 定理]
设 $D\subseteq\C$ 为单连通开集，$f$ 在 $D$ 内解析。则对 $D$ 内任意逐段光滑\textbf{闭曲线} $\gamma$（即 $\gamma(a)=\gamma(b)$），
\[
\oint_\gamma f(z)\,\d z = 0.
\]
\end{theorem}

\begin{theorem}[Cauchy 积分公式]
设 $D$ 为有界闭域，其边界 $\partial D$ 为逐段光滑正向闭曲线。$f$ 在 $D$ 内解析，在 $\partial D$ 上连续。则对任意 $z_0\in D$，
\[
f(z_0) = \frac{1}{2\pi i}\oint_{\partial D}\frac{f(z)}{z-z_0}\,\d z.
\]
\end{theorem}

\begin{theorem}[Cauchy 积分公式的推广]
在上述条件下，若在 $z_0$ 处 $f$ 有任意阶导数，则 $\forall n\in\N$，
\[
f^{(n)}(z_0) = \frac{n!}{2\pi i}\oint_{\partial D}\frac{f(z)}{(z-z_0)^{n+1}}\,\d z.
\]
\end{theorem}
\begin{example}
    利用积分 $\displaystyle\oint_{\abs{z}=1}\frac{\e^z}{z}\,\d z$，计算实积分
    \[\int_0^{2\pi}\e^{\cos\theta}\sin(\sin\theta)\,\d \theta,\qquad\int_0^{2\pi}\e^{\cos\theta}\cos(\sin\theta)\,\d \theta\]
\end{example}
\begin{solution}
    
    设 $z = e^{i\theta}$，则 $\d z = i e^{i\theta}\d\theta$，$\theta\in[0,2\pi]$。于是
    \[
    \oint_{\abs{z}=1}\frac{\e^z}{z}\,\d z = \int_0^{2\pi}\frac{\e^{e^{i\theta}}}{e^{i\theta}} i e^{i\theta}\d\theta = i\int_0^{2\pi}\e^{e^{i\theta}}\d\theta.
    \]
    由 Cauchy 积分公式可知
    \[
    \oint_{\abs{z}=1}\frac{\e^z}{z}\,\d z = 2\pi i\e^0 = 2\pi i.
    \]
    因此
    \[
    \int_0^{2\pi}\e^{e^{i\theta}}\d\theta = 2\pi.
    \]
    又有
    \[
    \int_0^{2\pi}\e^{e^{i\theta}}\d\theta = \int_0^{2\pi}\e^{\cos\theta}[\cos(\sin\theta)+i\sin(\sin\theta)]\,\d \theta.
    \]
    对比实部与虚部，得
    \[
    \int_0^{2\pi}\e^{\cos\theta}\cos(\sin\theta)\,\d \theta = 2\pi,\qquad
    \int_0^{2\pi}\e^{\cos\theta}\sin(\sin\theta)\,\d \theta = 0.
    \]
\end{solution}
\begin{example}
    计算积分 $\displaystyle\oint_C\frac{\e^{-z}\sin z}{z^2}\,\d z$，其中 $C:\abs{z-i}=2$
\end{example}
\begin{solution}
    由 Cauchy 积分公式的推广可知
    \[
    \oint_C\frac{\e^{-z}\sin z}{z^2}\,\d z = 2\pi i f'(0),
    \]
    其中 $f(z) = \e^{-z}\sin z$。于是
    \[
    f'(z) = -\e^{-z}\sin z + \e^{-z}\cos z = \e^{-z}(\cos z - \sin z),
    \]
    因此$f'(0) = 1$，从而
    \[
    \oint_C\frac{\e^{-z}\sin z}{z^2}\,\d z = 2\pi i.
    \]
\end{solution}
\ifx\allfiles\undefined
\end{document}
\fi